\documentclass[11pt]{article}
\usepackage[utf8]{inputenc}
\usepackage{amsmath,amssymb,amsthm}
\usepackage[margin=1in]{geometry}
\usepackage{enumitem}
\usepackage{hyperref}

\begin{document}
% ==================== PAGE 9 ====================
\noindent \textbf{Equivalence Class}

\noindent Let $*$ be a group action of a group $G$ on $A$
\[
* : G \times A \to A
\]
for $a, b \in A$, define a relation $\sim$ on $A$ as
\[
a \sim b \iff a = g * b, \quad \text{for some } g \in G
\]
then $\sim$ is an equivalence relation on $A$.

\medskip
\noindent \textbf{Proof:} \\
$\forall\, a \in A$
\[
a = e * a, \quad e \in G
\]
$\Rightarrow a \sim a$ \\
$\therefore \sim$ is reflexive.

\noindent Let $a, b \in A$ s.t. $a \sim b$ \\
$\Rightarrow \exists\, g \in G$ s.t.
\begin{align*}
a &= g * b  \\
\Rightarrow g^{-1} * a &= g^{-1} * (g * b)  \\
\Rightarrow g^{-1} * a &= b  \\
\Rightarrow b &\sim a
\end{align*}
$\therefore \sim$ is symmetric.

\noindent Let $a, b, c \in A$ s.t. $a \sim b$ and $b \sim c$ \\
$\Rightarrow \exists\, g_1, g_2 \in G$ s.t.
\[
a = g_1 * b \quad \text{and} \quad b = g_2 * c
\]
\[
a = g_1 * (g_2 * c)
\]
\[
a = (g_1 g_2) * c \quad (g_1 g_2 \in G)
\]
$\Rightarrow a \sim c$ \\
$\therefore \sim$ is transitive and equivalence relation.

\newpage
% ==================== PAGE 10 ====================
\noindent So, it will partition the set $A$ into it's equivalence class.

\noindent Let $a \in A$
\begin{align*}
\text{Cl}(a) &= \{b \in A \mid a \sim b\}  \\
&= \{b \in A \mid b = g * a, \; \forall\, g \in G\}  \\
&= \{g * a \mid g \in G\}
\end{align*}
\[
\text{Cl}(a) = G \cdot a \quad (\text{orbit of } a \text{ under } *)
\]
\[
\Rightarrow A = \bigcup_{a \in A} \text{Cl}(a)
\]
\[
\Rightarrow A = \bigcup_{a \in A} G \cdot a
\]

\bigskip
\noindent \textbf{Conjugacy class}

\noindent If $G$ acts on itself by conjugation, then
\[
G_a = \text{Normalizer of } a = N(a)
\]
\[
G \cdot a = \text{Cl}(a) = \{g * a \mid g \in G\}
\]
\[
= \{gag^{-1} \mid g \in G\} \quad (\text{conjugacy class of } a)
\]

\bigskip
\noindent \textbf{H.W} \\
Let $G$ be a group and $a, b \in G$. Then $a$ and $b$ are conjugate if $\exists\, c \in G$ such that $b = c a c^{-1}$. \\
Prove that $\sim$ is an equivalence relation on $G$ where $a \sim b$ when $a$ is conjugate to $b$.

\medskip
\noindent \textbf{Soln:} \\
Let $a \in G$, then
\[
a = e a e^{-1} \quad (\text{as } e \in G)
\]
$\Rightarrow a \sim a$ and $\sim$ is reflexive.

\newpage
% ==================== PAGE 11 ====================
\noindent Let $a, b \in G$ such that $a \sim b$, then $\exists\, c \in G$ s.t.
\begin{align*}
a &= c b c^{-1}  \\
c^{-1} a &= c^{-1} c b c^{-1}  \\
c^{-1} a c &= e b c^{-1} c  \\
c^{-1} a c &= b e  \\
c^{-1} a (c^{-1})^{-1} &= b
\end{align*}
$\Rightarrow b \sim a$ and $\sim$ is symmetric.

\noindent Let $a, b, c \in G$ s.t. $a \sim b$ and $b \sim c$ \\
then $\exists\, d_1, d_2 \in G$ s.t.
\begin{align*}
a &= d_1 b d_1^{-1}  \\
b &= d_2 c d_2^{-1}  \\
a &= d_1 d_2 c d_2^{-1} d_1^{-1} \quad \{d_2^{-1} d_1^{-1} = (d_1 d_2)^{-1}\}  \\
a &= (d_1 d_2) c (d_1 d_2)^{-1} \quad \{\text{as } d_1 d_2 \in G\}
\end{align*}
$\Rightarrow a \sim c$ and $\sim$ is transitive and hence equivalence.

\bigskip
\noindent \textbf{Note:} By orbit-stabilizer theorem $\exists$ a bijection between the set $\text{Cl}(a)$ and $G/N(a)$.
\[
G = \bigcup_{a \in G} \text{Cl}(a)
\]
1. If $G$ is a finite group then
\[
|\text{Cl}(a)| = \left|\frac{G}{N(a)}\right|
\]
\[
|G| = \sum_{a \in G} |\text{Cl}(a)|
\]

\bigskip
\noindent \textbf{To Prove:} $\text{Cl}(a) = \{a\} \iff a \in Z(G)$.

\newpage
% ==================== PAGE 12 ====================
\noindent \textbf{Proof:} \\
Let $\text{Cl}(a) = \{a\}$. \\
If $g \in G$ then
\begin{align*}
gag^{-1} \in \text{Cl}(a) &= \{a\}  \\
\Rightarrow gag^{-1} &= a  \\
\Rightarrow ga &= ag
\end{align*}
As $g$ is arbitrary so $a \in Z(G)$.

\noindent \textbf{Conversely,} let $a \in Z(G)$
\begin{align*}
\text{Cl}(a) &= \{gag^{-1} \mid g \in G\}  \\
&= \{agg^{-1} \mid g \in G\}  \\
\text{Cl}(a) &= \{a\}
\end{align*}

\bigskip
\noindent \textbf{To Prove:} $a \in Z(G) \iff N(a) = G$

\medskip
\noindent \textbf{Proof:} \\
Let $a \in Z(G)$
\begin{align*}
&\Rightarrow ag = ga \quad \forall\, g \in G  \\
&\Rightarrow g \in N(a) \quad \forall\, g \in G  \\
&\Rightarrow N(a) = G
\end{align*}
\textbf{Conversely,} let $N(a) = G$ \\
then $\forall\, g \in G, \; g \in N(a)$
\begin{align*}
&\Rightarrow ag = ga \quad \forall\, g \in G  \\
&\Rightarrow a \in Z(G)
\end{align*}

\bigskip
\noindent \textbf{Class Equation:}
\[
|G| = |Z(G)| + \sum_{a \notin Z(G)} |G / N(a)|
\]
\end{document}
